\documentclass[11pt,a4paper]{article}
\usepackage[T1]{fontenc}
\usepackage[utf8]{inputenc}
\usepackage[margin=2.3cm]{geometry}
\usepackage{lmodern}
\usepackage{microtype}
\usepackage{parskip}
\usepackage[table]{xcolor}
\usepackage{amsmath}
\usepackage{amssymb}
\usepackage{booktabs}
\usepackage{array}
\usepackage{enumitem}
\usepackage[numbers]{natbib}
\usepackage{hyperref}
\usepackage{titlesec}
\hypersetup{colorlinks=true, linkcolor=blue, urlcolor=blue, citecolor=blue}

\definecolor{chgcolor}{RGB}{0,80,180}
\definecolor{rccolor}{RGB}{0,110,40}

\newcommand{\rconcern}[3]{%
  \bigskip\noindent\rule{\linewidth}{0.4pt}\par\medskip
  \noindent\textbf{\textcolor{rccolor}{#1 --- #2:}}\par
  \smallskip
  \begin{quote}\itshape\textcolor{rccolor}{#3}\end{quote}%
}

\newcommand{\newtext}[1]{\textcolor{chgcolor}{``#1''}}

\newcommand{\msloc}[1]{%
  \smallskip\noindent\textit{Revised manuscript: #1}\par%
}

\titleformat{\section}{\large\bfseries}{}{0em}{}
\titlespacing*{\section}{0pt}{2.2em}{0.8em}

\begin{document}

% ------------------------------------------------------------------ Cover page
\noindent
Ms.\ Ref.\ No.: \textbf{KNOSYS-D-26-21490R1}\\[2pt]
Title: \textit{Effective Training Pressure Gates Recursive Knowledge
Degradation in LLMs: A Multi-Axis Dose-Response Study}\\[2pt]
Journal: \textit{Knowledge-Based Systems} (Elsevier)\\[2pt]
To: Dr.\ Hang Yu, Senior Editor\\[2pt]
Re: Response to Reviewers --- Production Submission\\[2pt]
Date: \today

\bigskip
\medskip\noindent\rule{\linewidth}{0.4pt}\medskip

\noindent
Dear Dr.\ Yu and Reviewers,

We thank you for the positive decision and for the careful reading of the
revised manuscript.
This letter accompanies the production files (LaTeX source and figures) and
responds point by point to the three questions raised by Reviewer~1 in the
acceptance decision.
Reviewer~2 expressed no further concerns; we are grateful for their
evaluation.

One textual change has been made to the manuscript: two sentences have been
added to Section~2.4 (Parameter-efficient fine-tuning and forgetting) to
address Reviewer~1's request for a comparison with sequential-task-adaptation
methods.
No other changes were made to the scientific content.

\medskip\noindent
Respectfully submitted on behalf of all authors,\\[6pt]
Julio Leite Azancort Neto\\
Corresponding Author

\medskip\noindent\rule{\linewidth}{0.4pt}\medskip

% ------------------------------------------------------------------ R2
\section*{Reviewer 2}

Reviewer~2 approved the manuscript without further concerns.
We are grateful for the thorough and constructive engagement throughout
both rounds of review.

% ------------------------------------------------------------------ R1
\section*{Reviewer 1}

\rconcern{Reviewer 1}{Question 1}{%
Can the authors report a targeted, blinded manual check of short and
verbose answers against the substring scorer, using the original fixed $K_0$
items, and explain whether any disagreements affect the key interpretations?}

The evaluation uses a bidirectional substring scorer: a generation is marked
correct if either (a)~the reference answer is a substring of the generated
response, or (b)~the generated response is a substring of the reference
answer.
Direction~(a) handles verbose generated responses that contain the correct
answer within a longer reply; direction~(b) handles concise generated
responses that partially match a longer reference alias.
The $K_0$ item set is fixed at generation~0 using this same scorer, so
the denominator is internally consistent with the numerator throughout all
generations.

During the revision process, the scoring pipeline was subjected to a full
audit that identified and corrected two implementation inconsistencies
(described in Section~3.3 of the revised manuscript).
After those corrections, we conducted a targeted manual review of a
stratified sample of $K_0$ items, covering both short reference answers
(fewer than five tokens) and verbose reference answers (more than fifteen
tokens).
For the short-answer stratum, the bidirectional check correctly handled all
sampled cases: short correct answers were captured by direction~(b) when
the generated response was also short, and by direction~(a) when the
generated response was verbose.
For the verbose-answer stratum, the check correctly retained items where
the full reference phrase appeared as a substring of the generation.
No disagreements were found in the sampled items that would change the sign
or statistical significance of any reported result.
We therefore conclude that the key interpretations are not affected by
answer-length variation.

\msloc{Section~3.3 (Evaluation protocol), revised manuscript.}

% -----------------------------------------------------------------
\rconcern{Reviewer 1}{Question 2}{%
Can the authors provide an immutable, internally consistent mapping from
each reported condition to its actual run, seed, precision, $K_0$ item
identifiers and generation-level counts, and regenerate the inferential
statistics from those records?}

The complete experimental record is publicly archived at:

\begin{center}
\url{https://doi.org/10.5281/zenodo.23145361}
\end{center}

The archive contains, for every reported condition, run directory, seed,
floating-point precision (NF4 or bf16), the $K_0$ item set constructed at
generation~0, and generation-level retention counts.
The mapping from manuscript condition labels to archive paths is provided
in the \texttt{README.md} at the repository root.

All inferential statistics reported in the manuscript (interaction estimates,
confidence intervals, sign tests, and TOST equivalence tests) were
regenerated directly from the per-seed \texttt{results\_all.json} files
archived at that DOI before this submission.
The regeneration scripts are included in the archive under
\texttt{scripts/audits\_diagnostics/}.

\msloc{Data Availability Statement and supplementary material,
revised manuscript.}

% -----------------------------------------------------------------
\rconcern{Reviewer 1}{Question 3}{%
The related-work discussion could further distinguish forgetting under
sequential task adaptation from degradation of a recursively replaced data
stream. A brief comparison with ``Adaptive tree-like neural network:
Overcoming catastrophic forgetting to classify streaming data with concept
drifts'' (Knowledge-Based Systems, 2024), focused on that distinction rather
than a direct LLM baseline, would help clarify the scope of the contribution.}

We have added two sentences to Section~2.4 (Parameter-efficient fine-tuning
and forgetting), immediately before the discussion of \citet{biderman2024}.
The added passage reads:

\begin{quote}
\newtext{Methods designed to overcome catastrophic forgetting under sequential
task adaptation, such as the adaptive tree-like neural network of
\citet{wen2024} which grows architecture branches for newly arriving concepts,
address a different failure mode from the one studied here: rather than
cross-task interference at discrete concept boundaries, the degradation we
observe arises from a recursively replaced data stream in which each
generation's synthetic outputs alter the distribution seen by the next.}
\end{quote}

\msloc{Section~2.4, paragraph beginning ``While classical continual
learning\ldots'', revised manuscript.}

\medskip
\noindent\textbf{Production files.} This submission includes (i)~LaTeX
source files (manuscript, bibliography, shared number definitions, and table
bodies), (ii)~figure files in PDF and PNG format at 300\,DPI, and
(iii)~updated Highlights.
The Highlights have been corrected to align with the manuscript: the bullet
previously reading ``arrests progressive loss'' has been replaced by ``A~50\%
synthetic exposure cut improved retention on both backbones tested,'' which
accurately reflects the distinct patterns observed on Qwen and Gemma~3.

\medskip\noindent\rule{\linewidth}{0.4pt}

\bibliographystyle{unsrtnat}
\bibliography{../../manuscript/cas-refs}

\end{document}
